\begin{codebox}
\codeline{\textcolor{cmtgray}{\#\ 描述逻辑公理与约束示例}}
\codeline{\textcolor{cmtgray}{\#\ Description\ Logic\ Axioms\ and\ Constraints}}
\codeline{\textcolor{cmtgray}{\#}}
\codeline{\textcolor{cmtgray}{\#\ 描述逻辑（Description\ Logic,\ DL）是OWL的理论基础}}
\codeline{\textcolor{cmtgray}{\#\ 提供严格的形式化语义和可判定的推理能力}}
\codeblank{}
\codeline{\textcolor{cmtgray}{\#\ ==============================}}
\codeline{\textcolor{cmtgray}{\#\ 1.\ 等价类定义\ (Equivalent\ Class)}}
\codeline{\textcolor{cmtgray}{\#\ ==============================}}
\codeline{\textcolor{cmtgray}{\#\ 用于定义类的充分必要条件，明确类的本质特征}}
\codeblank{}
\codeline{\textcolor{cmtgray}{\#\ 定义"合格产品"：}}
\codeline{QualifiedProduct\ \(\equiv\)\ Product\ AND\ (\(\forall\)\ inspection.hasResult(inspection,\ "Pass"))}
\codeline{\textcolor{cmtgray}{\#\ 即：合格产品是"属于Product类"且"所有检验的结果均为通过"的产品}}
\codeblank{}
\codeline{\textcolor{cmtgray}{\#\ ==============================}}
\codeline{\textcolor{cmtgray}{\#\ 2.\ 全称限制与存在限制}}
\codeline{\textcolor{cmtgray}{\#\ ==============================}}
\codeblank{}
\codeline{\textcolor{cmtgray}{\#\ 全称限制（Universal\ Restriction，\(\forall\)）：}}
\codeline{\textcolor{cmtgray}{\#\ \(\forall\)R.C\ 表示与该对象存在R关系的所有对象，都属于C类}}
\codeline{\(\forall\)R.C}
\codeblank{}
\codeline{\textcolor{cmtgray}{\#\ 存在限制（Existential\ Restriction，\(\exists\)）：}}
\codeline{\textcolor{cmtgray}{\#\ \(\exists\)R.C\ 表示该对象至少存在一个R关系的对象属于C类}}
\codeline{\(\exists\)R.C}
\codeblank{}
\codeline{\textcolor{cmtgray}{\#\ 使用构造器定义"高端设备"：}}
\codeline{HighEndEquipment\ \(\equiv\)\ Equipment\ \(\sqcap\)\ (price\ \(\ge\)\ 1000000)\ \(\sqcap\)\ (\(\forall\)hasFeature.AdvancedFeature)}
\codeline{\textcolor{cmtgray}{\#\ 即：高端设备是"属于设备类"、"价格大于等于100万"且"所有特征都是高级特征"的设备}}
\codeblank{}
\codeline{\textcolor{cmtgray}{\#\ ==============================}}
\codeline{\textcolor{cmtgray}{\#\ 3.\ 公理示例\ (Axiom\ Examples)}}
\codeline{\textcolor{cmtgray}{\#\ ==============================}}
\codeblank{}
\codeline{\textcolor{cmtgray}{\#\ 一致性公理示例（此两条公理相互矛盾）：}}
\codeline{\textcolor{cmtgray}{\#\ 公理1：所有设备都必须位于某个工作单元}}
\codeline{\(\forall\)x:\ Equipment(x)\ \(\rightarrow\)\ locatedIn(x,\ WorkCell)}
\codeline{\textcolor{cmtgray}{\#\ 公理2：存在不位于任何工作单元的设备}}
\codeline{\(\exists\)x:\ Equipment(x)\ \(\land\)\ \(\lnot\)(\(\exists\)w:\ locatedIn(x,\ w))}
\codeblank{}
\codeline{\textcolor{cmtgray}{\#\ ==============================}}
\codeline{\textcolor{cmtgray}{\#\ 4.\ 工程约束示例}}
\codeline{\textcolor{cmtgray}{\#\ ==============================}}
\codeblank{}
\codeline{\textcolor{cmtgray}{\#\ 推理优化：优先使用子类关系（\(\sqsubseteq\)）而非等价类（\(\equiv\)）}}
\codeline{\textcolor{cmtgray}{\#\ 减少等价类带来的推理负担}}
\codeblank{}
\codeline{\textcolor{cmtgray}{\#\ 只使用类层次和简单的存在约束}}
\codeline{\(\exists\)R.C}
\codeblank{}
\codeline{\textcolor{cmtgray}{\#\ 添加全称约束（推理复杂度增加）}}
\codeline{\(\forall\)R.C}
\codeblank{}
\codeline{\textcolor{cmtgray}{\#\ ==============================}}
\codeline{\textcolor{cmtgray}{\#\ 5.\ 概率本体（Probabilistic\ Ontology）}}
\codeline{\textcolor{cmtgray}{\#\ ==============================}}
\codeblank{}
\codeline{\textcolor{cmtgray}{\#\ 高功率设备需要维护的概率为85\%}}
\codeline{P(HighPowerEquipment\ \(\sqsubseteq\)\ NeedsMaintenance)\ =\ 85}
\codeline{\textcolor{cmtgray}{\#\ 含义：100台高功率设备中，约85台需要定期维护}}
\codeblank{}
\codeline{\textcolor{cmtgray}{\#\ 贝叶斯网络条件概率：}}
\codeline{P(Fault\ |\ HighTemp\ \(\land\)\ LongRunning)\ =\ 0.9\ \ \ \#\ 高温且长时间运行，故障概率90\%}
\codeline{P(Fault\ |\ HighTemp\ \(\land\)\ \(\lnot\)LongRunning)\ =\ 0.3\ \ \#\ 高温但未长时间运行，故障概率30\%}
\codeline{P(Fault\ |\ \(\lnot\)HighTemp\ \(\land\)\ LongRunning)\ =\ 0.1\ \ \#\ 未高温但长时间运行，故障概率10\%}
\codeline{P(Fault\ |\ \(\lnot\)HighTemp\ \(\land\)\ \(\lnot\)LongRunning)\ =\ 0.01\ \#\ 既未高温也未长时间运行，故障概率1\%}
\end{codebox}
